\documentclass{article}

\textwidth=6.5in
\textheight=8.5in
\voffset=-54pt
\oddsidemargin=0in
\evensidemargin=0in
\setlength{\parskip}{8pt}
\setlength{\parindent}{0pt}

\usepackage{workbook}

\usepackage[sols]{handout}

\begin{document}

\begin{center}
  \large\textsc{Problem Set 2 - Related Rates of Change}
  
  \pspace{12pt}
  \begin{problemonly}\large Name: \rule{5cm}{0.15mm}\\ \end{problemonly}
\end{center}

\begin{enumerate}

\item 

  A spherical balloon leaks air at a constant rate of 0.5
  $\text{cm}^3/\text{hour}$. We want to determine how quickly the radius is decreasing.
  \begin{enumerate}

  \item 
  \begin{problem}
      Write down a relationship between the radius $r$ of the balloon and its volume $V$.
  \end{problem}
  \pspace{\vfill}

  \begin{solution}
  The formula for the volume of a sphere in terms of its radius is $V=\frac{4}{3}\pi r^3$.
  \end{solution}

  \item
  \begin{problem}
      Using Leibniz notation, write down the rate we are given and the rate we want to determine.
  \end{problem}
  \pspace{\vfill}

  \begin{solution}
      We are given the rate of change of the volume, which is $\diff{V}{t}$ in Leibniz notation. Since we are told that the volume is decreasing, we know that $\diff{V}{t} = -0.5$ cubic centimeters per hour.

      We want to determine the rate of change of the radius, $\diff{r}{t}$. 
  \end{solution}

  \item
\begin{problem}[\stretch{3}]
      Find how fast the radius is decreasing when the diameter of the balloon is 30 cm.  Include units in your answer.
    \end{problem}

    \begin{solution} 

        \textbf{Rate we know:} $\diff{V}{t} = -0.5 \text{ cm}^3/\text{hr}$ 

        \textbf{Rate we want to know:} $\diff{r}{t}[r=15\text{ cm}]$

      \textbf{Equation relating our variables:} $V =
        \frac{4}{3} \pi r^3$

      Once we have our relating equation, we differentiate it.
        Since we are looking for $\frac{dr}{dt}$, we should
        differentiate with respect to $t$.
      \begin{align}
        \frac{d}{dt} V &= \frac{d}{dt} \!\left( \frac{4}{3} \pi r^3
        \right) \nonumber \\
        \frac{dV}{dt} &= 4 \pi r^2
        \frac{dr}{dt}
\\
        \shortintertext{Plugging in what we know,}
        -0.5 &= 4 \pi (15)^2 \frac{dr}{dt}, \nonumber 
      \end{align}
      so $\frac{dr}{dt} = \frac{-0.5}{4 \pi (15)^2} = - \frac{1}{1800
        \pi}$ cm/hour.  That is, the radius is decreasing at an
        instantaneous rate of $\frac{1}{1800 \pi}$ cm/hour when the diameter of the balloon is 30 cm.
    \end{solution}

  \item
    \begin{problem}[\stretch{1}]
      How quickly is the radius decreasing when the diameter of the
      balloon is 20 cm?
    \end{problem}

    \begin{solution}
      This is essentially the same as {{{{\#1}}(a)}},
      except that we want the rate at the moment when $r = 10$.  We can
      still use equation (1), but now we plug in $r
      = 10$ to get $-0.5 = 4 \pi (10)^2 \frac{dr}{dt}$, which gives
      $\frac{dr}{dt} = - \frac{1}{800 \pi}$.  So, the radius is
      decreasing at an instantaneous rate of $\frac{1}{800 \pi}$
        cm/hour when the diameter is 20 cm.
    \end{solution}

  % \item
  %   \begin{problem}[\nstretch{1}]
  %     Based on geometry, you think the radius should decrease more or less quickly as the balloon gets smaller?  Do your answers agree with
  %     this?
  %   \end{problem}

  %   \begin{solution}
  %     If the volume of air is decreasing at a constant rate,
  %     we would expect the radius to decrease \textbf{more quickly} as the
  %     balloon gets smaller (for example, losing 1 cm$^3$ of air when the
  %     balloon is big will barely affect the radius, but losing 1 cm$^3$ of
  %     air when the balloon is tiny will have a huge effect on the radius).
  %     This agrees with our answers.
  %   \end{solution}

  \end{enumerate}


\pspace{\newpage}

\item 
  \begin{problem}
    Ana, Jelena, and Novak are working on the following problem.
    \begin{quote} % Stewart 4.1, #25
      Gravel is being dumped from a conveyor belt at a rate of 30
      ft$^3$/min into a cylindrical container.  How
      fast is the height of the gravel in the container increasing when the gravel has reached 10 ft high? 
    \end{quote}
    They've decided to use the variables $V = $ volume of gravel, $h = $
    height of gravel, and $t$ = time.  Now, they're trying to decide on a
    strategy.
    \begin{itemize}[nolistsep]
    \item Ana says, ``We're looking for $\frac{dh}{dt}$, so we should come up with a formula for
      $h$ as a function of $t$ and then differentiate.''
    \item Jelena says, ``We're looking for $\frac{dh}{dt}$, but we know
      $\frac{dV}{dt}$, so we should first relate $h$ and $V$ and then
      differentiate with respect to $t$.''
    \item Novak says, ``We're looking for $\frac{dh}{dV}$, so we should
      first relate $h$ and $V$ and then differentiate with
      respect to $V$.''
    \end{itemize}
    Whose strategy makes the most sense? Explain briefly.
  \end{problem}
  \pspace{\vfill}

  \begin{solution}
    Jelena's strategy makes the most sense. We are looking for $\frac{dh}{dt}$ and we are told $\frac{dV}{dt}$. We won't be able to find a direct relationship between $h$ and $t$, so we need to relate $h$ and $V$.
  \end{solution}

\item
\begin{grading}
4 points. 2 points for each correct answer. Deduct a point for each incorrect answer.
\end{grading}

\begin{problem}
  João and Tanya are out taking walks.  Let $x$ be the distance between
  them at time $t$.  If João is currently 1 mile south of Tanya, in which
  of the following situations is $\frac{dx}{dt}$ negative right now?
  List all.
  \begin{enumerate}[label=\protect\maybecircle{\Alph*.}]
  \plainitem João is walking south at a rate of 3 mph, while Tanya is walking
    south at a rate of 2.5 mph.
  \circleitem João is walking south at a rate of 2.5 mph, while Tanya is
    walking south at a rate of 3 mph.
  \plainitem João is walking north at a rate of 3 mph, while Tanya is
    walking north at a rate of 3.5 mph.
  \circleitem João is walking north at a rate of 3 mph, while Tanya is
    walking south at a rate of 3.5 mph.
  \plainitem João is walking south at a rate of 3 mph, while Tanya is
    walking north at a rate of 3.5 mph.
  \end{enumerate}
  \emph{Hint:} There are two correct answers. A picture will help you!
\end{problem}
\pspace{\vfill}

  \begin{solution} % Jason Bland
    $\frac{dx}{dt}$ is negative when $x$ is decreasing over time.
     This happens when João and Tanya are getting closer to each
    other, so the correct choices are B and D.
  \end{solution}

\pspace{\newpage}


\item 
\begin{problem}
    An ant is crawling along a path in the shape of an ellipse centered at the origin described by the equation $4x^2+9y^2=25$. 
\end{problem}

\begin{enumerate}
    \item
    \begin{problem}
        When the ant is at the point $(2,1)$, $dy/dt=\frac{1}{3}$. Find $dx/dt$. Is the ant moving clockwise or counterclockwise around the ellipse at this moment?
    \end{problem}
    \pspace{\vfill}

    \begin{solution}
        Since we know $dy/dt$ and want to find $dx/dt$, we can take the relationship between $x$ and $y$ and differentiate both sides with respect to $t$.
        \begin{align*}
            \diff*{(4x^2+9y^2)}{t} &= \diff*{25}{t} \\
            8x\diff{x}{t} + 18y\diff{y}{t} &= 0\\
        \intertext{This gives us a relationship between $dx/dt$ and $dy/dt$. We can substitute the values we have for $x$, $y$, and $dy/dt$ to solve for $dx/dt$.}
            8(2)\diff{x}{t} + 18(1)(\tfrac{1}{3}) &= 0 \\
            \diff{x}{t} &= -\frac{3}{8}
        \end{align*}
        Since $dy/dt$ is positive and $dx/dt$ is negative, the ant is moving up and to the left. This means that, since the ellipse is centered at the origin, when the ant is at the point $(2,1)$, it is moving counterclockwise around the ellipse.
    \end{solution}

        \item
    \begin{problem}
        When the ant is at the point $(-2,1)$, $dy/dt=3$. Find $dx/dt$. Is the ant moving clockwise or counterclockwise around the ellipse at this moment?
    \end{problem}
    \pspace{\vfill}

    \begin{solution}
        In the last part, we found a relationship between $dx/dt$, $dy/dt$, $x$, and $y$ that we can use to answer this part by substituting the values of $dy/dt$, $x$, and $y$ and solving for $dx/dt$.
        \begin{align*}
            8x\diff{x}{t} + 18y\diff{y}{t} &= 0\\
            8(-2)\diff{x}{t} + 18(1)(3) &= 0\\
            \diff{x}{t} &= \frac{27}{8}
        \end{align*}
        Since $dy/dt$ is positive and $dx/dt$ is positive, the ant is moving up and to the right. This means that, since the ellipse is centered at the origin, when the ant is at the point $(-2,1)$, it is moving clockwise around the ellipse.
    \end{solution}
\end{enumerate}

\pspace{\newpage}

\item 

\begin{problem}
    A lamp is mounted at the top of a 15-ft-tall pole. A bear that is 6 ft tall stands 30 ft from the pole. How long is the bear's shadow? (Drawing a picture may help.)
\end{problem}
\pspace{\stretch{1.5}}

\begin{solution}
    Let's draw a diagram of the situation.

     \begin{center}
        \begin{tikzpicture}

% Define coordinates
\coordinate (O) at (0,0); % Base of the street lamp
\coordinate (P) at (3,0); % Foot of the bear's shadow
\coordinate (Q) at (5,0); % Total shadow
\coordinate (L) at (0,1.5); % Top of the street lamp
\coordinate (M) at (3,0.6); % Top of the bear

% Draw street lamp
\draw[very thick] (O) -- node[midway, above, rotate=90] {15 ft} (L);

% Draw woman
\draw[very thick] (P) -- node[midway, left] {6 ft} (M);

% Label shadow lengths and heights
\draw[very thick] (O) -- node[midway, below] {30 ft} (P);
\draw[very thick] (P) -- node[midway, below] {$x$} (Q);
\draw[very thick] (L) -- (Q);

\end{tikzpicture}
    \end{center}
We can use similar triangles to set up a proportion and solve for $x$.
\begin{align*}
    \frac{x}{6} &= \frac{x+30}{15} \\
    \intertext{We can multiply both sides by 30 to clear the fractions.}
    5x &= 2(x+30) \\ 
    5x &= 2x+60\\
    3x &= 60 \\
    x &= 20
\end{align*}
The bear's shadow is 20 feet long.
\end{solution}

\item
\begin{problem}
    A conical water cup measures 2 inches in diameter and 4 inches tall. The formula for the volume of a cone with radius $r$ and height $h$ is $V=\frac{1}{3}\pi r^2 h$.
\end{problem}
    
    \begin{enumerate}
        \item 
        \begin{problem}
        How much water is in the cup when it is filled to a height of 2 inches? (Drawing a picture may help. \emph{Hint}: You need to use similar triangles.)
        \end{problem}
        \pspace{\vfill}

        \begin{solution}
            Let's draw a diagram.
            \begin{center}
        \begin{tikzpicture}
          \def\R{1}
          \def\H{4}
          \def\r{0.5}
          \def\h{2}
          \def\aspect{15}
          \draw [fill=white] (-\R,\H) -- (\R,\H) -- (0,0) -- cycle;
          \draw [fill=white] (0, \H) circle [x radius=\R, y
          radius={\R/\aspect}];
          \draw [fill=blue!40!white,opacity=1] (-\r, \h) -- (\r, \h) --
          (0,0) -- cycle; 
          \draw [fill=blue!20!white,opacity=1, draw opacity=0] (0, \h) ellipse
          [x radius=\r, 
          y radius={\r/\aspect}];
          \draw (-\r, \h) arc [x radius=\r, y radius={\r/\aspect}, start
          angle=180, end angle=360];      
          \draw[dashed] (\r, \h) arc [x radius=\r, y radius={\r/\aspect},
          start angle=0, end angle=180];
          \draw[blue, very thick] (0, 0) -- (0, \H) --
          node[above, black]{1 in}
          (\R, \H) -- cycle;      
          \draw[blue, very thick] (0, \h) -- node[above,
          black]{$r$} (\r, \h); 
          \draw[|-|, xshift=-24pt] (-\R, 0) -- node[left]{4 in} (-\R,
          \H);
          \draw[|-|, xshift=-12pt] (-\r, 0) -- node[left]{2 in} (-\r,
          \h);
        \end{tikzpicture}
      \end{center}    
      When the cup is filled with water, the water forms a smaller, similar cone inside. Since the height of the water is half the height of the cup, the radius of the surface of the water will be half the radius of the cup, so $r=0.5$ inches. (We could have also set up a proportion to find $r$.)

      Plugging $r=0.5$ inches and $h=2$ inches into the formula for the volume of a cone gives us $V=\tfrac{1}{3}\pi(0.5)^2(2)=\frac{1}{6}\pi$ cubic inches.
        \end{solution}
        
        \item
        \begin{problem}
What fraction of the total volume of the cup is this?
    \footnote{If you find this strange and want to know more, search for the video ``Cones are MESSED UP" by Numberphile.}
    \end{problem}
    \pspace{\stretch{0.6}}

    \begin{solution}
        The total volume of the cup is $\tfrac{1}{3}\pi(1)^2(4)=\frac{4}{3}\pi$ cubic inches. This means that the amount of water in the cup when it is filled halfway to the top is
        \[\frac{\frac{1}{6}\pi}{\frac{4}{3}\pi} = \frac{1}{6} \cdot \frac{3}{4} =  \frac{1}{8}\]
        of the total volume of the cup.
    \end{solution}
    
    \end{enumerate}

\pspace{\newpage}

\item 
  \note{Finding trigonometric values for $30^\circ, 45^\circ, 60^\circ$.} 

  The two right triangles shown are examples of the two families of ``special right triangles.'' The length of the shortest side of both triangles as shown is 1 unit long.  
  \begin{center}
    \begin{tabular}{ccc}
      an isosceles right triangle & & half of an equilateral triangle \\ \\
      \begin{tikzpicture}[x=2.5cm, y=2.5cm]
        \draw (0, 0) -- node[below]{1} (1, 0) -- (1, 1) -- cycle;
        \draw (1,0) rectangle +(-0.1, +0.1);
      \end{tikzpicture}
      &
      \hspace{0.25in} &
      \begin{tikzpicture}[x=1.5cm, y=1.5cm]
        \draw (0, 0) -- node[below]{1} (1, 0) -- (1, {sqrt(3)}) -- cycle;
        \draw[dashed] (1, {sqrt(3)}) -- (2, 0) -- (1, 0);
        \draw (1,0) rectangle +(-0.15, +0.15);
      \end{tikzpicture} \\      
    \end{tabular}  
  \end{center}
  \begin{enumerate}
  \item
    \begin{problem}
      Find and label the lengths of all other sides of the triangles, as
      well as the size of the angles in each triangle.
    \end{problem}

    \begin{solution}
      Let's start with the right isosceles triangle.  ``Isosceles'' means
      the two legs have the same length, so both legs have length 1.  By the
      Pythagorean Theorem, the hypotenuse therefore has length $\sqrt{2}$.
      The angles of this triangle are $\ang{45}$, $\ang{45}$, and
      $\ang{90}$.  Thus, here is a completely labeled triangle:
      \begin{center}
        \begin{tikzpicture}[x=2.5cm, y=2.5cm]
          \draw (0, 0) -- node[below]{1} (1, 0) -- node[right]{1} (1, 1)
          -- node[anchor=south east]{$\sqrt{2}$} (0, 0);
          \draw (8pt, 0) arc[radius=8pt, start angle=0, end angle=45]
          node[anchor=180]{$\ang{45}$};
          \draw ($(1, 1) - (0, 8pt)$) arc[radius=8pt, start angle=270, end
          angle=225] node[anchor=100]{$\ang{45}$};
        \end{tikzpicture}
      \end{center}
      For the equilateral triangle, since half of one side has length 1,
      the sides of the equilateral triangle are all 2.  Therefore, the
      left half is a right triangle with base 1 and hypotenuse 2; by the
      Pythagorean Theorem, its height is $\sqrt{3}$.  The angles of the
      equilateral triangle are all $\ang{60}$, and half of this is
      $\ang{30}$:
      \begin{center}
        \begin{tikzpicture}[x=1.5cm, y=1.5cm]
          \draw (0, 0) -- node[below]{1} (1, 0) -- node[pos=0.3, right,
          inner sep=0.5pt]{$\sqrt{3}$}
          (1, {sqrt(3)}) -- node[anchor=south east]{$2$} (0, 0);
          \draw[dashed] (1, {sqrt(3)}) -- (2, 0) -- (1, 0);
          \draw (8pt, 0) arc[radius=8pt, start angle=0, end angle=60]
          node[anchor=180]{$\ang{60}$};
          \draw ($(1, {sqrt(3)}) - (0, 8pt)$) arc[radius=8pt, start
          angle=270, end angle=240] node[anchor=95]{$\ang{30}$};
        \end{tikzpicture}        
      \end{center}
    \end{solution}

\item 
  \begin{problem}
    It's often useful to work with right triangles that have a hypotenuse that is 1 unit long. If we scale down both of the triangles from before so that the hypotenuse is 1 unit long,  Find and label the
    side lengths and angles of the smaller triangles.
     \begin{center}
    \begin{tabular}{ccc}
      an isosceles right triangle & & half of an equilateral triangle \\ \\
      \begin{tikzpicture}[x=2.5cm, y=2.5cm]
        \draw (0, 0) --  (1, 0) -- (1, 1) -- node[midway, above left] {1} (0,0);
        \draw (1,0) rectangle +(-0.1, +0.1);
      \end{tikzpicture}
      &
      \hspace{0.25in} &
      \begin{tikzpicture}[x=1.5cm, y=1.5cm]
        \draw (0, 0) -- (1, 0) -- (1, {sqrt(3)}) -- node[midway, above left] {1} (0,0);
        \draw (1,0) rectangle +(-0.15, +0.15);
      \end{tikzpicture} \\      
    \end{tabular}  
  \end{center}
  \end{problem}

  \begin{solution}
    The first triangle is similar to the right isosceles triangle
    in {{\#1}}.  Since this triangle has hypotenuse
    1 and the one in {{\#1}} had hypotenuse
    $\sqrt{2}$, the sides of this one are all $\frac{1}{\sqrt{2}}$ as big
    as the ones in {{\#1}}; therefore, the two legs
    both have length $\frac{1}{\sqrt{2}}$.

    The second triangle is similar to the ``half of an equilateral
    triangle'' in {{\#1}}; the sides of this one are
    $\frac{1}{2}$ as long as the ones in {{\#1}}:
    \begin{center}
      \begin{tikzpicture}[x=2.5cm, y=2.5cm]
        \draw (0, 0) -- node[below]{$\frac{1}{\sqrt{2}}$} (1, 0) --
        node[right]{$\frac{1}{\sqrt{2}}$} (1, 1) -- node[anchor=south
        east]{$1$} (0, 0); %
        \draw (8pt, 0) arc[radius=8pt, start angle=0, end angle=45]
        node[anchor=180]{$\ang{45}$}; %
        \draw ($(1, 1) - (0, 8pt)$) arc[radius=8pt, start angle=270, end
        angle=225] node[anchor=100]{$\ang{45}$}; %
      \end{tikzpicture}
      \hspace{0.25in}
      \begin{tikzpicture}[x=1.5cm, y=1.5cm]
        \draw (0, 0) -- node[below]{$\frac{1}{2}$} (1, 0) -- node[pos=0.3,
        right, 
        inner sep=0.5pt]{$\frac{\sqrt{3}}{2}$} (1, {sqrt(3)}) --
        node[anchor=south east]{1} (0, 0); %
        \draw[dashed] (1, {sqrt(3)}) -- (2, 0) -- (1, 0); %
        \draw (8pt, 0) arc[radius=8pt, start angle=0, end angle=60]
        node[anchor=180]{$\ang{60}$}; %
        \draw ($(1, {sqrt(3)}) - (0, 8pt)$) arc[radius=8pt, start
        angle=270, end angle=240] node[anchor=95]{$\ang{30}$}; %
      \end{tikzpicture}        
    \end{center}    
  \end{solution}

\end{enumerate}

\item
\begin{problem}
    Sarah measures the shortest side of a 30-60-90 triangle to be $7$ cm long. How long is the hypotenuse?
\end{problem}
\pspace{\vfill}

\begin{solution}
    In a 30-60-90 triangle, the hypotenuse is twice as long as the shortest side. So the hypotenuse of Sarah's triangle is 14 cm long.
\end{solution}

\end{enumerate}

\psreflection{
  \emph{Reminders:}
  \begin{itemize}[topsep=0pt]
   \item Complete Math Ma Skills Check Practice 1 and turn it in by \textbf{Thursday, January 30th at 5:30pm.}
  \item The Math Ma Skills Check
  is \textbf{Thursday, February 6th 5:30-6:30pm}.  If you have an approved conflict, you must fill out the Out of Sequence Exam form (found on Canvas) by Friday, January 30th at 7pm to arrange an earlier time to take the Skills Check.
  \end{itemize}

    Recommended reading for next class: examples 17.11 and 17.12 on pp. 552-554 of Gottlieb.
}

\end{document}